% !TEX root = ../main.tex
\section{Scale, Scope, and What Is Actually Tunable}
\label{sec:scale}

The architectural choices of Sections~\ref{sec:governed-rag} and~\ref{sec:bounded}
were presented as though the model were fixed. It usually is not. The question an
architect faces early is whether to consume a large general-purpose model as a
service or to operate a smaller model over a defined domain, and the two options are
routinely compared on accuracy. This section argues that accuracy is close to
irrelevant to the comparison, identifies the two variables that do govern it, and
declines a tempting conclusion that does not follow.

\subsection{What scale does, and what it cannot do}

Scale improves accuracy. That is not in dispute and nothing here argues otherwise.
What matters is what an accuracy improvement licenses, and by
Proposition~\ref{prop:independence} it licenses nothing about the grounding rate. The
two quantities are not linked by any entailment, so a model that
scores better tells you it scores better.

There is a stronger point available, and it follows from
Section~\ref{sec:pretraining} rather than from any claim about how well these systems
work. The objective \eqref{eq:objective} contains no term that is a function of
truth. Accuracy gains from scale are therefore gains in modelling the training
distribution, which correlates with truth to the degree the corpus is truthful. They
are not convergence toward truth as an optimisation target, because no term in the
loss distinguishes a true continuation from a plausible false one where the two
diverge. Scale sharpens the estimate of the distribution; it does not introduce a
distinction the objective never drew.

This is worth stating carefully, because it is easy to overstate in either
direction. The claim is not that a larger model is not more often right---it is. The
claim is that being more often right on a distribution is a different property from
standing in an auditable relation to evidence, and that no amount of the first
produces the second. The mechanism that would produce grounding is absent from the
architecture, per Section~\ref{sec:missing}, and adding parameters does not add a
mechanism.

\subsection{Reproducibility is a property of topology, not of size}

Invariant I14 is where the comparison is usually miscast. Reproducibility does not
scale, in either direction, because it is not a function of the model at all. It is a
function of who operates the inference.

Section~\ref{sec:determinism} listed the conditions: pinned model versions,
controlled batching or single-request inference, deterministic kernel selection,
recorded decoding parameters, statelessness. Every one of those is a property of the
serving stack. A parameter count appears in none of them.

What is true, and is the defensible core of the common intuition, is that the
\emph{typical consumption pattern} for large models makes those conditions harder to
obtain. A shared multi-tenant endpoint batches your request with other traffic, so
batch composition is outside your control; sparse expert routing under capacity
limits makes that batch composition affect your output; and the served version may
change without notice. None of that follows from the model being large. It follows
from the model being operated by someone else.

The distinction is not pedantic, because it changes what an architect should do about
it. ``Prefer a small model'' is not actionable advice and is sometimes wrong.
``Obtain version pinning, deterministic serving, and stateless requests---by
self-hosting, by contract, or by choosing a deployment that offers them'' is
actionable, and it is available for open-weight models of any size.

\subsection{Corpus scope changes what can be measured}
\label{sec:corpus-scope}

Here the comparison is genuinely asymmetric, and in a way that has nothing to do with
parameter count and everything to do with what the evidential base contains.

Two of this paper's measurements require establishing a \emph{negative}: that the
evidence does not support a claim. The abstention items of
Section~\ref{sec:abstention} need it by construction---an item whose correct
behaviour is silence is only well formed if you can establish that the answer is
absent from what the system can draw on. And the ungrounded assertion rate
\eqref{eq:rate} needs it for the parametric half of $E$, which is why
Section~\ref{sec:measurability} concluded that half is largely unmeasurable.

Negative claims about a bounded, enumerated, versioned corpus are tractable. You can
determine that a controlled document set does not address a question, by search over
an enumerated and versioned index rather than by inference about what a model
absorbed. For a corpus of any size that is work rather than inspection, and it is
tractable in the way any completeness check over a known population is tractable. Negative claims about a web-scale training corpus are not
tractable with current attribution methods, and the gap is one of kind rather than
maturity: the question is what a model did \emph{not} absorb, and attribution
techniques estimate influence of what it did.

A related asymmetry follows from Section~\ref{sec:floor}. The statistical floor on
ungrounded assertion is highest for facts that are rare in the training
distribution---and \emph{rare} is only definable relative to a known corpus. For a
domain corpus you can characterise coverage and therefore estimate where the floor
sits. For a frontier training set you cannot, so the floor is present, consequential,
and unquantified.

So bounding the corpus does not merely improve accuracy on the domain. It changes
which properties of the system are \emph{measurable at all}, which is a different and
more important kind of gain. It is the same move as
Section~\ref{sec:bounded}, applied to the evidence rather than the request:
enumerability is what makes negative statements available, and negative statements
are what auditing requires.

\subsection{What the dials actually buy}

Collecting the choices available to an architect against the invariants they affect:

\begin{table}[ht]
  \centering
  \small
  \begin{tabular}{@{}>{\raggedright\arraybackslash}p{4.0cm}>{\raggedright\arraybackslash}p{2.4cm}>{\raggedright\arraybackslash}p{6.4cm}@{}}
    \toprule
    Choice & Invariants & Effect on auditability \\
    \midrule
    Typed request interface
      & I1, I7, I10
      & Large. Supplies the canonical form and enforces context binding \\
    \addlinespace
    Controlled, versioned corpus
      & I8, I9, I11
      & Large. Precondition for every evidence condition \\
    \addlinespace
    Bounded, enumerable corpus
      & enables I13, I17 testing
      & Large. Makes negative evidence claims and abstention items possible \\
    \addlinespace
    Schema-constrained output plus external gate
      & I12, I13, I17
      & Large. Atomic errors; abstention enforced rather than expressed \\
    \addlinespace
    Operational control of inference
      & I14, I7
      & Moderate to large. Version pinning, batching, statelessness \\
    \addlinespace
    Claim decomposition
      & I12
      & Moderate. Claim-level rather than response-level linkage \\
    \addlinespace
    \textbf{Model parameter count}
      & \textbf{none directly}
      & \textbf{No direct effect.} Affects accuracy, which by
        Prop.~\ref{prop:independence} bounds nothing. One indirect path matters---see
        below \\
    \bottomrule
  \end{tabular}
  \caption{Architectural choices ranked by effect on auditability. Model size is the
    dial most often debated and the one with the least direct effect.}
  \label{tab:dials}
\end{table}

The ordering is the practical conclusion of this section, and it is somewhat
counterintuitive: the variable that dominates procurement discussion has the smallest
\emph{direct} effect on whether the resulting system can be audited.

One indirect path is real and should not be buried in a table. Where an in-scope model
performs the context resolution or request construction of
Section~\ref{sec:relocation}---translating free text into a typed request---its
capability bears directly on I10, and a misparse there produces the most dangerous
failure in the architecture: a complete, internally consistent record answering the
wrong question. So model capability does matter, at exactly one point, and the
conclusion to draw is not that a larger model fixes it. It is that the parse must be
made \emph{reviewable}, per Section~\ref{sec:relocation}, because a control that depends
on model capability is a control whose failure rate is unknown. Better capability
narrows that failure rate without bounding it; confirmation bounds it.

\subsection{The conclusion that does not follow}
\label{sec:not-follow}

It is tempting to conclude that small domain models can be governed and large
general-purpose models cannot. That does not follow, and it is worth saying why,
because the error would misdirect real spending.

Every condition in Table~\ref{tab:invariants-system} is a property of the system
surrounding the model. None mentions parameter count, training corpus size, or
provenance of the weights. A large model placed inside a governed architecture---typed
requests, controlled corpus, external evidence gate, defeater search, audit record,
pinned version---satisfies those conditions exactly as a small model would, because the
conditions are about the architecture. Conversely, a small model consumed through a
shared endpoint with no version pinning and free-form input satisfies almost none of
them.

The two genuine axes identified above are \emph{operational control of the inference
stack} and \emph{enumerability of the evidential base}. Both correlate with model
size in current practice, and neither is caused by it. Self-hosted open weights give
operational control at any scale. A retrieval corpus can be bounded and versioned
whatever model reads it.

The claim I would defend is therefore narrower than the intuition and more useful
than it: \textbf{you cannot govern inference you neither operate nor contract for, and
you cannot make negative evidence claims about a corpus you cannot enumerate.}

The contractual half deserves emphasis, because enterprises govern third-party
processors routinely and the instruments are familiar. What would otherwise be
configuration becomes a term, and the terms are nameable: version pinning with a
defined support window; advance notice of model and serving changes; a deterministic or
single-tenant serving option; recorded decoding parameters; retention of request and
response for a stated period; audit rights over the serving configuration; and
restrictions on sub-processing. An organisation that cannot obtain those has not
discovered a property of large models---it has discovered the terms it accepted. The
practical conclusion is a procurement checklist rather than a prohibition, and it is
the one place in this paper where the binding control is legal rather than technical.

One asymmetry does survive, and it is worth stating because it is where the intuition
is right. Scale is offered as a solution to accuracy, and accuracy is not the
constraint. So an organisation that responds to a grounding problem by buying a larger
model has bought an improvement it cannot use for the purpose it needs, and has
usually given up operational control to obtain it. That is not an argument against
large models. It is an argument against treating model capability as the lever when
the binding constraint is architectural.
